\chapter{Geltungsbereich und Grundbegriffe}

Dieser Leitfaden ergänzt Band~7 um zentrale Definitionen,
Charakterisierungen und Folgerungen für \emph{surjektive Funktionen
zwischen Mengen}. Grundlage sind die Zermelo--Fraenkel-Mengenlehre
ohne Auswahlaxiom (ZF) und der Funktionsbegriff aus Band~5.
Beweise und Beweisskizzen entfallen.

\section{Gegenstand und Abgrenzung}

Aufgenommen werden Kriterien für Surjektivität, ihre besonderen
mengentheoretischen Folgen und Konstruktionen surjektiver Funktionen.
Allgemeine Funktionsgesetze werden vorausgesetzt; ihre Notation
wird kurz festgelegt. Endlichkeit ist keine allgemeine Voraussetzung.

Die Theorien injektiver und bijektiver Funktionen, Äquivalenz- und
Ordnungsrelationen, Quotienten, Mächtigkeiten und zusätzliche
algebraische oder topologische Strukturen bleiben ihren eigenen
Bänden vorbehalten. Mengeninklusionen und Gleichheiten von
Funktionswerten dienen hier nur der Formulierung der Aussagen.

Das Auswahlaxiom wird weder vorausgesetzt noch verneint.
Die allgemeine Existenz von Rechtsinversen sowie die Surjektivität
beliebiger unendlicher Produkte werden nicht behauptet;
Auswahlprinzipien behandelt Band~9.
Bedingungen wie Nichtleere oder eine Überdeckung stehen jeweils
im betreffenden Satz.

Die Zusammenstellung ist eine Auswahl bekannter Hauptresultate,
keine vollständige Liste aller gefundenen oder ableitbaren Aussagen.
Aufbau und Gestaltung folgen den Leitfäden zu Band~3, Band~4,
Band~5, Band~28 und Band~37.
\section{Surjektivität und Hilfsnotation}

\begin{definition}[Surjektive Funktion]
\label{b07guide:def:surjection}
Eine Funktion $f\colon A\to B$ heißt \emph{surjektiv auf $B$}, wenn
\[
 \forall b\in B\;\exists a\in A\quad f(a)=b.
\]
Wir schreiben dafür $f\colon A\twoheadrightarrow B$.
Mit $B^A=\{f\mid f\colon A\to B\}$ sei
\[
 \operatorname{Sur}(A,B)
   =\{f\in B^A\mid f\text{ ist surjektiv auf }B\}.
\]
\end{definition}

Wie im Hauptband ist eine Funktion ihr Graph, also eine Menge
geordneter Paare. Eine zulässige Zielmenge ist kein zusätzlicher
Bestandteil dieses Graphen. Die Aussage, daß ein Graph
\emph{auf eine angegebene Menge} surjektiv ist, hängt jedoch von
dieser Zielmenge ab. Funktionswerte werden nur im angegebenen
Definitionsbereich verwendet.

\begin{definition}[Bild, Urbild und Faser]
Für $f\colon A\to B$, $U\subseteq A$, $V\subseteq B$ und $b\in B$ sei
\[
 \begin{aligned}
 f[U]&=\{f(a)\mid a\in U\},\\
 f^{-1}[V]&=\{a\in A\mid f(a)\in V\},\\
 \operatorname{Fib}_f(b)&=f^{-1}[\{b\}].
 \end{aligned}
\]
$f^{-1}[V]$ bezeichnet eine Urbildmenge, keine Umkehrfunktion.
$\mathcal P(A)$ bezeichnet die Potenzmenge von $A$.
\end{definition}

\begin{definition}[Einschränkung, Komposition und Identität]
Für $f\colon A\to B$ und $U\subseteq A$ sei
\[
 f|_U=\{(a,f(a))\mid a\in U\}.
\]
Für $g\colon B\to C$ ist $(g\circ f)(a)=g(f(a))$ für $a\in A$.
Die Identität auf $A$ lautet
$\operatorname{id}_A\colon A\to A$, $a\mapsto a$.
\end{definition}

\section{Wegweiser durch die Resultate}

\begin{center}
\begin{tabularx}{\textwidth}{@{}p{.30\textwidth}X@{}}
\toprule
Fragestellung & Fundstelle\\
\midrule
Surjektivität erkennen &
Bild-, Faser- und Teilmengenkriterien sowie leere Randfälle
in Kapitel~\ref{b07guide:ch:criteria}.\\
Zielmengen prüfen &
Rückgewinnung, Urbildtests und Quantorenübertragung
in Kapitel~\ref{b07guide:ch:images}.\\
Bereiche verändern &
Zielmengenwechsel, Einschränkung und Zusammenfügen
in Kapitel~\ref{b07guide:ch:restrictions}.\\
Kompositionen prüfen &
Komposition, Rechtskürzen und eindeutige Faktorisierung
in Kapitel~\ref{b07guide:ch:composition}.\\
Neue Surjektionen bilden &
Projektionen, Produkte, markierte Vereinigungen und
Potenzmengen in Kapitel~\ref{b07guide:ch:constructions}.\\
Unmöglichkeit erkennen &
Cantors Satz am Ende von Kapitel~\ref{b07guide:ch:constructions}.\\
\bottomrule
\end{tabularx}
\end{center}

Die Nummern mit Präfix $7H$ gehören ausschließlich zu diesem
Leitfaden. Die Nummerierung des vollständigen Fachbands bleibt
unverändert.